\begin{frame}{One network, many affine maps}
  \small
  \setlength{\abovedisplayskip}{6pt}\setlength{\belowdisplayskip}{6pt}
  \begin{columns}[c,onlytextwidth]
    \begin{column}{0.51\textwidth}
      Each ReLU either passes its input or sets it to zero:
      \[
        \vc{h}=\vc{D}(\vc{x})(\vc{W}\vc{x}+\vc{b}),
        \qquad D_{ii}(\vc{x})\in\{0,1\}.
      \]
      The diagonal entries of $\vc{D}$ indicate which units are active. With an affine output,
      \begin{align*}
        f(\vc{x})&=\vc{w}^{\top}\vc{D}(\vc{x})(\vc{W}\vc{x}+\vc{b})+c\\
                 &=\vc{a}(\vc{x})^{\top}\vc{x}+\beta(\vc{x}).
      \end{align*}
      Within each region, the gates stay fixed: the network is \myalert{exactly affine}.

      \medskip
      This holds at any depth. Training learns both the regions and their affine maps.
    \end{column}
    \begin{column}{0.46\textwidth}
      \centering
      $f(\vc{x})=\operatorname{ReLU}(x_1)+\operatorname{ReLU}(x_2-x_1)$\par
      \smallskip
      % Original NN figure, 2026-10-01. Exact activation partition for
% f(x) = ReLU(x_1) + ReLU(x_2-x_1). See verify_geometry.py.
\begin{tikzpicture}[x=1.14cm,y=1.14cm]
  \path[use as bounding box] (-2.35,-2.45) rectangle (2.55,2.50);
  \fill[drawgray_l] (-2,-2)--(0,0)--(0,-2)--cycle;
  \fill[drawblue_l!75] (-2,-2)--(-2,2)--(0,2)--(0,0)--cycle;
  \fill[draworange_l!75] (0,-2)--(2,-2)--(2,2)--(0,0)--cycle;
  \fill[drawgreen_l!85] (0,0)--(2,2)--(0,2)--cycle;
  \draw[drawgray!45] (-2,-2) rectangle (2,2);
  \draw[->,drawgray!70] (-2.15,0)--(2.30,0) node[right,text=black] {$x_1$};
  \draw[->,drawgray!70] (0,-2.15)--(0,2.30) node[above,text=black] {$x_2$};
  \draw[drawgray,thick,densely dashed] (0,-2)--(0,2);
  \draw[drawgray,thick,densely dashed] (-2,-2)--(2,2);
  \node[anchor=west,font=\scriptsize,text=drawgray] at (1.40,2.25) {$x_2=x_1$};
  \node[below right,font=\scriptsize,text=drawgray] at (0.03,-0.03) {$0$};
  \node[align=center] at (-0.60,-1.46) {$f=0$\\[-1pt]{\scriptsize $(0,0)$}};
  \node[align=center] at (-1.03,0.86) {$f=x_2-x_1$\\[-1pt]{\scriptsize $(0,1)$}};
  \node[align=center] at (1.08,-0.89) {$f=x_1$\\[-1pt]{\scriptsize $(1,0)$}};
  \node[align=center] at (0.57,1.47) {$f=x_2$\\[-1pt]{\scriptsize $(1,1)$}};
  \node[font=\scriptsize,text=drawgray] at (0,-2.36) {Labels: output formula and gate pattern $(d_1,d_2)$};
\end{tikzpicture}

    \end{column}
  \end{columns}
  \vfill
  {\scriptsize\color{drawgray}For ReLU hidden layers, before any output sigmoid or softmax.
  \href{https://arxiv.org/abs/1402.1869}{Mont\'ufar et al. (2014), Linear regions of deep neural networks}.}
\end{frame}

\begin{frame}{Depth through composition}
  \small
  \setlength{\abovedisplayskip}{5pt}\setlength{\belowdisplayskip}{5pt}
  Three ReLUs make a triangular function. Apply the same function repeatedly:
  \[
    T(x)=2\operatorname{ReLU}(x)-4\operatorname{ReLU}\!\left(x-\tfrac12\right)
         +2\operatorname{ReLU}(x-1),\qquad x\in[0,1].
  \]
  \begin{center}
    % Original NN plot of the ReLU triangle construction used by Telgarsky (2016).
% T^k has 2^k affine pieces on [0,1]. Vertices are exact dyadic fractions.
\begin{tikzpicture}[x=3.4cm,y=1.70cm]
  \path[use as bounding box] (-0.10,-0.47) rectangle (3.80,1.49);
  \foreach \shift/\pieces/\heading in {0/2/{T(x)},1.36/4/{T(T(x))},2.72/8/{T(T(T(x)))}} {
    \begin{scope}[shift={(\shift,0)}]
      \foreach \j in {1,...,\pieces} {
        \pgfmathsetmacro{\leftedge}{(\j-1)/\pieces}
        \pgfmathsetmacro{\rightedge}{\j/\pieces}
        \ifodd\j
          \fill[drawblue_l!55] (\leftedge,0) rectangle (\rightedge,1);
        \else
          \fill[drawgreen_l!65] (\leftedge,0) rectangle (\rightedge,1);
        \fi
      }
      \draw[drawgray!35,densely dotted] (0,1)--(1,1);
      \draw[->,drawgray] (-0.04,0)--(1.06,0) node[right,text=black] {$x$};
      \draw[->,drawgray] (0,-0.04)--(0,1.13);
      \node[below left,font=\scriptsize] at (0,0) {$0$};
      \node[below,font=\scriptsize] at (1,0) {$1$};
      \node[left,font=\scriptsize] at (0,1) {$1$};
      \foreach \j in {1,...,\pieces} {
        \pgfmathsetmacro{\xa}{(\j-1)/\pieces}
        \pgfmathsetmacro{\xb}{\j/\pieces}
        \ifodd\j
          \draw[darkcerulean,line width=1.25pt] (\xa,0)--(\xb,1);
        \else
          \draw[darkcerulean,line width=1.25pt] (\xa,1)--(\xb,0);
        \fi
      }
      \node at (0.5,1.35) {$\heading$};
      \node[font=\small] at (0.5,-0.34) {\textbf{\pieces} affine pieces};
    \end{scope}
  }
  \draw[-{Stealth},drawgray,thick] (1.10,0.54)--(1.24,0.54);
  \draw[-{Stealth},drawgray,thick] (2.46,0.54)--(2.60,0.54);
\end{tikzpicture}

  \end{center}
  Each composition splits every piece in two: $k$ blocks give $2^k$ pieces using $3k$ ReLU units.

  \smallskip
  A single hidden layer with $m$ ReLUs has at most $m+1$ pieces in one dimension.
  \myalert{Depth can make a representation much more compact.}

  \smallskip
  Without nonlinearities, all layers collapse to one affine map.
  \vfill
  {\scriptsize\color{drawgray}An expressivity example, not a guarantee about training or generalization.
  \href{https://proceedings.mlr.press/v49/telgarsky16.html}{Telgarsky (2016), Benefits of depth in neural networks}.}
\end{frame}
